% !TeX root = ../navier-stokes-manuscript.tex
\subsection{A spatial column inside every temporal response}\label{sub:TheViscosityDrivenSpatialAscent}

Before estimating anything, fix the algebraic point at which the two directions
meet.  Begin with \(V_0=u\), whose critical height first exposed the spatial
movement.  For \(s\geq0\), set
\[
  E_s(t):=\frac12\norm{\Lambda^su(t)}_2^2
\]
and retain the complete transport--pressure pairing
\[
  \mathcal P_s(t)
  :=-
  \operatorname{Re}
  \pair{\Lambda^su}
  {\Lambda^s\bigl((u\cdot\nabla)u+\nabla p\bigr)}_{L^2}.
\]
The momentum equation paired at that height is
\begin{equation}\label{eq:BodyFullSpatialEnergyBalance}
  E_s'=\mathcal P_s-2\nu E_{s+1}.
\end{equation}
At \(s=\tfrac12\),
\begin{equation}\label{eq:BodyFirstCriticalRow}
  H'=\mathcal P_{1/2}-2\nu E_{3/2}.
\end{equation}
Differentiate this same row.  The first two steps are
\begin{align}
  H''
  &=\partial_t\mathcal P_{1/2}
  -2\nu\mathcal P_{3/2}
  +(2\nu)^2E_{5/2},
  \label{eq:BodySecondCriticalRow}\\
  H'''
  &=\partial_t^2\mathcal P_{1/2}
  -2\nu\partial_t\mathcal P_{3/2}
  +(2\nu)^2\mathcal P_{5/2}
  -(2\nu)^3E_{7/2}.
  \label{eq:BodyThirdCriticalRow}
\end{align}
Each differentiation exposes one higher spatial energy, and each lower
transport--pressure transfer survives beside it.

At finite order \(n\), the complete identity is
\begin{equation}\label{eq:BodyCriticalHeightSpatialAscent}
  H^{(n)}
  =(-2\nu)^nE_{n+1/2}
  +\sum_{j=0}^{n-1}
  (-2\nu)^j
  \partial_t^{\,n-1-j}\mathcal P_{j+1/2},
  \qquad n\geq1.
\end{equation}
Its induction from \eqref{eq:BodyFullSpatialEnergyBalance} is carried out in
\Cref{prop:CriticalHeightSpatialAscent}.  Because \(\nu>0\), the row can be
solved for its top coordinate:
\begin{equation}\label{eq:BodyCompletedCriticalRow}
  E_{n+1/2}
  =\frac{
    H^{(n)}-
    \displaystyle\sum_{j=0}^{n-1}
    (-2\nu)^j\partial_t^{\,n-1-j}\mathcal P_{j+1/2}
  }{(-2\nu)^n}.
\end{equation}
The division has not erased the transfer sum.  It has isolated the exact
collection that must be controlled before the exposed spatial energy is
controlled.  This is where the hoped-for reversal becomes conditional:
viscosity has put a higher energy in the equation, but writing the equation has
not bounded the nonlinear chain on its right.

The height derivative remains a contraction of the velocity responses:
\begin{equation}\label{eq:BodyHeightDerivativeVelocityContraction}
  H^{(n)}
  =\frac12\sum_{a=0}^{n}\binom{n}{a}
  \operatorname{Re}
  \pair{\Lambda^{1/2}V_a}
       {\Lambda^{1/2}V_{n-a}}_{L^2}.
\end{equation}
The three readings at order \(n\) must remain distinct.  \(V_n\) is the
Eulerian response of the velocity, \(H^{(n)}\) is a quadratic contraction of
the temporal responses, and \(E_{n+1/2}\) is the highest spatial energy
exposed after the full transfer chain has been retained.

\begin{center}
\captionsetup{hypcap=false}
\captionof{table}{The temporal response and the spatial height exposed at the same order.}
\label{tab:CompletedCriticalHeightRows}
\small
\renewcommand{\arraystretch}{1.3}
\begin{tabular}{@{}c c c c@{}}
\toprule
order & velocity response & critical response & exposed height \\
\midrule
\(0\) & \(V_0=u\) & \(H=E_{1/2}\) & \(E_{1/2}\) \\
\(1\) & \(V_1=u_t\) & \(H'\) & \(E_{3/2}\) \\
\(2\) & \(V_2=u_{tt}\) & \(H''\) & \(E_{5/2}\) \\
\(3\) & \(V_3=u_{ttt}\) & \(H'''\) & \(E_{7/2}\) \\
\(n\) & \(V_n=\partial_t^nu\) & \(H^{(n)}\) & \(E_{n+1/2}\) \\
\bottomrule
\end{tabular}
\end{center}

Now shift the base from \(u\) to any fixed response \(V_r\).  The Laplacian
still acts linearly on that same response, and every split of the transport
derivative remains in the binomial sum.  The question is no longer whether a
higher spatial coordinate appears.  It is where the equation actually pays
for that coordinate on the history already in hand.
